\documentclass[conference]{IEEEtran}
\IEEEoverridecommandlockouts

\usepackage{cite}
\usepackage{amsmath,amssymb,amsfonts}
\usepackage{algorithmic}
\usepackage{graphicx}
\usepackage{textcomp}
\usepackage{xcolor}
\usepackage{booktabs}
\usepackage{array}

\usepackage{url}
\usepackage{microtype}

\newcommand{\CN}{\mathcal{CN}}
\newcommand{\vect}[1]{\boldsymbol{#1}}
\newcommand{\mat}[1]{\mathbf{#1}}
\newcommand{\herm}{\mathrm{H}}
\newcommand{\diag}{\operatorname{diag}}
\newcommand{\argmin}{\operatorname*{arg\,min}}
\def\BibTeX{{\rm B\kern-.05em{\sc i\kern-.025em b}\kern-.08em
    T\kern-.1667em\lower.7ex\hbox{E}\kern-.125emX}}

\begin{document}
\bstctlcite{BSTcontrol}

\title{Fixed-Noise Sparse Bayesian Learning for Passive Intermodulation Cancellation}

\author{\IEEEauthorblockN{Anonymous Authors}}

\maketitle

\begin{abstract}
We evaluate fixed-noise sparse Bayesian learning (SBL) and target-correlation screening for complex Volterra-type passive intermodulation (PIM) cancellation. Six validation-tuned methods are compared on independent fitting, validation and test records, with ten repetitions, two nonlinear sources and four noise conditions. Screening at $\eta=0.01$ retains every candidate; higher thresholds can reduce the dictionary, whereas relevance pruning produces compact final models. On the clean cubic source, ridge achieves $58.02\pm0.87$ dB and SBL $51.96\pm0.65$ dB. At 35 dB PIM-to-noise ratio, their cancellation is approximately 35 dB, with 126 and 7.8 active coefficients on average. For tanh with a cubic--quintic dictionary at the same noise level, screen-plus-SBL achieves $28.33\pm0.46$ dB with 24 coefficients on average, versus $28.31\pm0.50$ dB with 4788 for ridge. Sequential resource measurements and support-overlap diagnostics distinguish coefficient economy from fitting cost and basis stability. Results characterize a fixed synthetic protocol and finite tuning grids; neither universal superiority nor calibrated predictive uncertainty is claimed.
\end{abstract}
\begin{IEEEkeywords}
passive intermodulation, sparse Bayesian learning, Volterra model, correlation screening, nonlinear system identification
\end{IEEEkeywords}

\section{Introduction}\label{sec:intro}
Passive intermodulation (PIM) arises when strong transmitted signals excite weak nonlinearities in nominally passive structures. In frequency-division-duplex radios, products generated in connectors, antennas or external objects can overlap the receiver band \cite{Lui1990PIM,Lampu2021AirInduced,Lampu2023Algorithms,Ahmmed2023Overview}. Digital cancellation reconstructs these products from known transmit samples and subtracts the replica from the received waveform.

Volterra and memory-polynomial models express nonlinear interactions through a regression linear in its coefficients, but their size grows rapidly with nonlinear order and memory \cite{Schetzen1980Volterra,Morgan2006GMP,Waheed2020PIM}. Physical and cascaded structures reduce this cost and have been evaluated on measured PIM \cite{Lampu2023Algorithms,Lampu2023Cascade}. Liu et al. combine multi-band Wiener--Hammerstein structure, FIR filters and spline nonlinearities for PIM reconstruction \cite{Liu2024MultiBand}; their structural reduction differs from selecting columns in a fixed Volterra dictionary.

Sparse Bayesian learning (SBL) is an established basis-selection method \cite{Tipping2001SBL,Wipf2004SBL}. Crespo-Cadenas et al. select frequency-domain regressors for OFDM power-amplifier modeling \cite{Crespo2025SparseRF}; that work concerns amplifier distortion rather than the receiver-filtered PIM target considered here. Wei combines SBL with BIC and penalized error-to-signal criteria to refine nonlinear dynamical models from small data \cite{Wei2025Identification}. Our experiment instead compares a one-shot target-correlation screen followed by fixed-noise SBL with unscreened SBL, ridge, LASSO and OMP under one controlled protocol. Applying SBL or placing a correlation test before it is not, by itself, a new inference principle. Published cancellation numbers from different signals and targets are not treated as direct baselines.

Our contribution is a controlled evaluation of whether correlation screening adds value beyond relevance pruning in a frequency-selected PIM dictionary. We examine symmetry-reduced cubic and cubic--quintic dictionaries, distinguish nonlinear-order mismatch from finite-memory truncation, and report sizes before screening, after screening and after pruning. The completed benchmark contains 720 validation-selection/test tasks and matched screening controls. Historical single-realization results are retained only for reproduction and convergence diagnostics.

\section{PIM Regression Model}
\label{sec:model}
\subsection{Two-Carrier Complex-Baseband Model}
Let $x_{\ell}[n]$ and $x_h[n]$ denote the known low- and high-frequency complex-baseband transmit-side reference components, and let $y[n]$ denote the extracted PIM observation in the receiver band corresponding to $2f_{\ell}-f_h$. A behavioral canceller represents the target as
\begin{equation}
  \mat y=\mat X_{\mathrm{full}}\mat w+\mat e,
  \label{eq:linear-model}
\end{equation}
where $\mat X_{\mathrm{full}}$ is assembled from delayed nonlinear monomials, $\mat w$ is the coefficient vector, and $\mat e$ contains modeling error and, in the noisy experiments, additive noise. Although the mapping from transmit samples to PIM is nonlinear, \eqref{eq:linear-model} is linear in $\mat w$. %In the supplied notebooks, the monomials are formed from the available transmit-side references rather than from oracle signals at the nonlinear source.
In the implementation, the monomials are formed from the available transmit-side references rather than from oracle signals at the nonlinear source.
Each raw monomial $\phi_i[n]$ is passed through the same known receiver low-pass filter used to isolate the selected PIM band,
\begin{equation}
 x_i[n]=(h_{\mathrm{rx}}*\phi_i)[n],
 \label{eq:filtered-column}
\end{equation}
and the filtered sequences form the columns of $\mat X_{\mathrm{full}}$. The causal memory terms therefore also represent the pre-nonlinearity linear-path dependence seen from the transmitter references.

For a causal delay set $\mathcal D$ of cardinality $M$, the cubic candidates that fall at $2f_{\ell}-f_h$ are
\begin{equation}
 \phi^{(3)}_{a,b,c}[n]
 =x_{\ell}[n-a]x_{\ell}[n-b]\overline{x_h[n-c]},
 \quad a\leq b,\quad a,b,c\in\mathcal D.
 \label{eq:cubic}
\end{equation}
The ordering $a\leq b$ removes permutation duplicates, so the number of distinct cubic columns is
\begin{equation}
 N_3(M)=\binom{M+1}{2}M=\frac{M^2(M+1)}{2}.
 \label{eq:n3}
\end{equation}
Thus $N_3=6,40,126,$ and $288$ for $M=2,4,6,$ and $8$.

The combined dictionary adds two fifth-order families that land in the same PIM band:
\begin{align}
 \phi^{(5a)}[n]&=x_{\ell,a}x_{\ell,b}x_{\ell,c}
                 \overline{x_{\ell,d}}\,\overline{x_{h,e}},
                 \quad a\leq b\leq c, \label{eq:im5a}\\
 \phi^{(5b)}[n]&=x_{\ell,a}x_{\ell,b}\overline{x_{h,c}}
                 \,\overline{x_{h,d}}x_{h,e},
                 \quad a\leq b,\ c\leq d, \label{eq:im5b}
\end{align}
where $x_{q,a}$ abbreviates $x_q[n-a]$. The symmetry-reduced candidate count is
\begin{equation}
\begin{split}
 N_{3+5}(M)={}&\binom{M+1}{2}M
 +\binom{M+2}{3}M^2\\
 &+\binom{M+1}{2}^{\!2}M,
\end{split}
\label{eq:n35}
\end{equation}
which gives 40, 760, 4788, and 18336 columns for $M=2,4,6,$ and $8$, respectively. These values describe the full symmetry-reduced candidate pool before correlation screening and Bayesian pruning; they are used only in the complexity and sparsity comparisons below. For clarity, $M$ denotes the number of causal taps per signal, rather than the half-width of a two-sided delay grid. The ordering restrictions remove algebraically duplicate columns without changing the span of the candidate model.

\subsection{Dense Ridge-Regularized Least-Squares Reference}
%For every memory/order configuration, the notebook baseline fits the full symmetry-reduced dictionary and is not subjected to correlation screening.
For every memory/order configuration, the reference estimator fits the full symmetry-reduced dictionary and is not subjected to correlation screening.
Let $\mat X_f\in\mathbb C^{N_f\times P}$ and $\mat y_f$ denote the fitting samples. The implementation forms
\begin{align}
 \mat R&=\frac{\mat X_f^{\herm}\mat X_f}{N_f},\qquad
 \mat r=\frac{\mat X_f^{\herm}\mat y_f}{N_f}, \label{eq:ls-stats}\\
 \lambda_X&=\rho\,\frac{\operatorname{tr}(\mat R)}{P},\qquad
 \rho>0, \label{eq:ls-lambda}
\end{align}
and computes
\begin{equation}
\begin{split}
 \widehat{\mat w}_{\mathrm{LS}}
 &=\left(\mat R+\lambda_X\mat I\right)^{-1}\mat r\\
 &=\argmin_{\mat w}\left\{
 \frac{1}{N_f}\|\mat y_f-\mat X_f\mat w\|_2^2
 +\lambda_X\|\mat w\|_2^2\right\}.
\end{split}
\label{eq:ls}
\end{equation}
The historical sweeps used $\rho=10^{-4}$; the revised protocol selects $\rho$ on validation data. The actual penalty $\lambda_X$ is scaled by the mean diagonal level of the fitting covariance. %When $P>N_f$, the fifth-order notebooks evaluate the algebraically equivalent dual form, avoiding a $P\times P$ solve without changing the estimator. 
When $P>N_f$, the fifth-order implementation evaluates the algebraically equivalent dual form, avoiding a $P\times P$ solve without changing the estimator.
Ridge stabilization improves conditioning and shrinks all coefficients, but it does not select an exactly sparse support. Throughout the tables and supplied plot legends, the short label ``LS'' denotes this dense ridge-regularized baseline.

\section{Correlation-Pruned Fixed-Noise SBL}
\label{sec:method}
\subsection{Normalized Correlation Screening}
Let $\mat y_f$ and $\mat x_{i,f}$ denote the fitting-sample target and the fitting-sample restriction of column $i$ of $\mat X_{\mathrm{full}}$. The target and each nonzero column are normalized as
\begin{equation}
 \bar{\mat y}_f=\frac{\mat y_f}{\|\mat y_f\|_2},\qquad
 \bar{\mat x}_{i,f}=\frac{\mat x_{i,f}}{\|\mat x_{i,f}\|_2}.
 \label{eq:normalize}
\end{equation}
Each candidate receives the score
\begin{equation}
 R_i=\left|\bar{\mat y}_f^{\herm}\bar{\mat x}_{i,f}\right|.
 \label{eq:corr}
\end{equation}
Given a threshold $\eta$, the one-shot screened set is
\begin{equation}
 \mathcal S_{\eta}=\{i:R_i>\eta\}.
 \label{eq:screen}
\end{equation}
The implementation retains the normalized screened matrix $\bar{\mat X}_f=[\bar{\mat x}_{i,f}]_{i\in\mathcal S_{\eta}}$ for the subsequent SBL iterations, rather than reverting to the original column scales. In the reported runs, optional forced top-$K$ retention is disabled and at least one candidate must remain; if the threshold rejects every nonzero column, the single highest-scoring column is retained. Normalization makes the screen and Bayesian updates invariant to individual column scale. Whether screening actually removes candidates must be measured; it is not implied by the presence of the screening step.

\subsection{Fixed-Noise Bayesian Model}
The normalized screened regression is assigned a complex Gaussian likelihood and independent coefficient prior,
\begin{align}
 p(\bar{\mat y}_f\mid\mat v,\beta)
   &=\CN(\bar{\mat X}_f\mat v,\beta^{-1}\mat I), \\
 p(\mat v\mid\vect\alpha)
   &=\CN(\mat 0,\mat A^{-1}),\qquad
     \mat A=\diag(\alpha_1,\ldots,\alpha_K).
 \label{eq:priors}
\end{align}
Here $\alpha_i$ is a coefficient precision: a large value means a small prior variance and therefore an irrelevant basis function. The likelihood precision $\beta$ is held constant inside one SBL run. Because the fitting target and dictionary columns are normalized, $\beta$ is an algorithmic precision in normalized coordinates; it is not estimated from the physical AWGN level. A single run therefore corresponds to one chosen value of $\beta$; the historical sweeps repeat the same solver at $\beta\in\{10^2,10^4,10^6\}$, but $\beta$ is never updated by the SBL recursion. No random samples are added to clean data merely to execute the algorithm. This external precision sweep is distinct from the experiment in which AWGN is physically added to the observation.

For fixed $\vect\alpha$ and $\beta$, the posterior covariance and mean in normalized coordinates are
\begin{align}
 \mat\Sigma&=\left(\beta\bar{\mat X}_f^{\herm}\bar{\mat X}_f+\mat A\right)^{-1},
 \label{eq:sigma}\\
 \vect\mu&=\beta\mat\Sigma\bar{\mat X}_f^{\herm}\bar{\mat y}_f.
 \label{eq:mu}
\end{align}
The effective degree of freedom and precision update are
\begin{align}
 \gamma_i&=1-\alpha_i\operatorname{Re}\{\Sigma_{ii}\}, \label{eq:gamma}\\
 \alpha_i^{\mathrm{new}}&=\frac{\gamma_i}{|\mu_i|^2}. \label{eq:alpha-update}
\end{align}
Numerically, the implementation clips $\gamma_i$ to $[\varepsilon,1]$, adds $\varepsilon$ to the denominator in \eqref{eq:alpha-update}, and clips the resulting precision to $[10^{-12},10^{16}]$, with $\varepsilon=10^{-12}$. Non-finite values or values above $\alpha_{\max}$ are removed; if every candidate would be removed, the candidate with the smallest updated precision is retained. Because $\vect\mu$ refers to normalized columns and a normalized target, the surviving coefficients are returned to the original dictionary scale. If $c_y=\|\mat y_f\|_2$ and $c_i=\|\mat x_{i,f}\|_2$ are the fitting-sample norms, then
\begin{equation}
 \widehat w_i=\frac{c_y}{c_i}\mu_i,\qquad
 \widehat{\mat p}=\mat X_{\mathrm{full}}[:,\mathcal A]
 \widehat{\mat w}_{\mathcal A},\qquad
 \mat r=\mat y-\widehat{\mat p}.
 \label{eq:estimate}
\end{equation}

\begin{figure}[!t]
\centering
\begin{minipage}{0.98\columnwidth}
\small
\begin{algorithmic}[1]
\REQUIRE $\mat y,\mat X_{\mathrm{full}},\eta,\beta,\alpha_{\max},T_{\max}$
\STATE Compute fitting-sample norms; normalize $\mat y_f$ and all nonzero columns
\STATE Compute $R_i=|\bar{\mat y}_f^{\herm}\bar{\mat x}_{i,f}|$ and retain $R_i>\eta$; if empty, retain the largest $R_i$
\STATE Form $\bar{\mat X}_f$ from the normalized retained columns
\STATE Initialize $\vect\alpha\leftarrow\mat 1$ on the screened dictionary
\FOR{$t=1,\ldots,T_{\max}$}
  \STATE Evaluate $\mat\Sigma$ and $\vect\mu$ from \eqref{eq:sigma}--\eqref{eq:mu}
  \STATE $\gamma_i\leftarrow\operatorname{clip}(1-\alpha_i\operatorname{Re}\{\Sigma_{ii}\})$
  \STATE $\alpha_i\leftarrow\operatorname{clip}(\gamma_i/(|\mu_i|^2+\varepsilon))$
  \STATE Remove non-finite or over-threshold precisions; keep one if needed and reset the previous mean after pruning
  \IF{no pruning occurred and the relative change in $\vect\mu$ is below tolerance}
    \STATE \textbf{break}
  \ENDIF
\ENDFOR
\STATE Recompute the posterior mean using the final active precisions
\STATE Rescale the surviving coefficients according to \eqref{eq:estimate}
\RETURN $\widehat{\mat p}$ and active basis indices
\end{algorithmic}
\end{minipage}
\caption{%Correlation-pruned fixed-noise SBL as implemented in the reported notebook sweeps. Posterior inference is performed on normalized variables, numerical safeguards are applied to the precision update, and the final coefficients are rescaled before PIM synthesis.
Correlation-pruned fixed-noise SBL with posterior inference on normalized variables, numerical safeguards in the precision update, and coefficient rescaling before PIM synthesis.}
\label{fig:algorithm}
\end{figure}


Algorithm~\ref{fig:algorithm} separates one-shot screening from iterative relevance pruning. In the historical experiments the first stage retains every candidate at $\eta=0.01$; sparsity is therefore attributable to the second stage. The revised solver caches the Gram matrix on each active set, uses Hermitian solves, and recomputes the final posterior after the last precision update. This last operation is an explicit numerical variant; the historical convergence diagnostic below preserves the original update order.

The two conjugated high-band factors in \eqref{eq:im5b} carry indices $c,d$, hence $c\leq d$; $e$ indexes the non-conjugated factor. The implementation stores the non-conjugated index before the ordered conjugated pair, an equivalent relabeling verified against direct monomial construction. Permutation duplicates are removed without assuming numerical linear independence of the remaining columns.

\section{Experimental Rationale and Setup}\label{sec:setup}
\subsection{Synthetic Signal and Receiver Chain}
Each record contains $N_s=10240$ complex samples at $f_s=256$ MHz. A single sequence $u[n]=a[n]+jb[n]$ is drawn using NumPy RandomState, with independent $a,b\sim U[0,1)$. It is neither centered nor normalized in power. A 2049-tap Hamming low-pass prototype $h_0$ is generated by \texttt{firwin(2049,0.05)} with unit DC gain. The transmit components are
\begin{equation}
 z_q[n]=\sum_{k=0}^{2048}h_0[k]e^{j2\pi f_qk/f_s}u[n-k].
 \label{eq:tx-generation}
\end{equation}
Here $f_\ell=-30$ MHz and $f_h=30$ MHz. Thus the carrier separation is 60 MHz, the prototype cutoff is 6.4 MHz, and its nominal two-sided bandwidth is 12.8 MHz; these are filter specifications, not measured occupied bandwidth. Both components are obtained from one white sequence by modulating the filter coefficients, rather than by modulating one already filtered waveform. They are not independently generated OFDM transmissions.

The composite $z_\ell+z_h$ passes through a causal 11-tap complex FIR,
\begin{equation}
 s[n]=\sum_{m=0}^{10}h[m](z_\ell[n-m]+z_h[n-m]).
 \label{eq:linear-path}
\end{equation}
Table~\ref{tab:path} specifies the fixed synthetic path. There are no negative delays, centered convolution or group-delay compensation. The available regression references are $x_q[n]=z_q[n]e^{-j2\pi f_qn/f_s}$, before this path; the path therefore contributes effective memory to the target.
\begin{table}[!t]
\caption{Coefficients of the causal synthetic linear path.}\label{tab:path}
\centering\scriptsize
\begin{tabular}{rr@{\qquad}rr}\toprule
$m$ & $h[m]$ & $m$ & $h[m]$\\\midrule
0 & $0.186-0.023j$ & 6 & $1.633-1.725j$\\
1 & $-0.098+0.117j$ & 7 & $-1.415+1.365j$\\
2 & $0.291-0.338j$ & 8 & $0.760-0.668j$\\
3 & $-0.607+0.628j$ & 9 & $-0.256+0.211j$\\
4 & $1.131-1.348j$ & 10 & $0.045-0.037j$\\
5 & $-1.090+0.993j$ & & \\\bottomrule
\end{tabular}
\end{table}

The nonlinear output is shifted by $+90$ MHz and passed through a causal 2049-tap Hamming receiver low-pass, \texttt{firwin(2049,0.2)}, with cutoff 25.6 MHz. This selects the product centered at $2f_\ell-f_h=-90$ MHz. Every dictionary monomial undergoes the same receiver filter. The first $D=2\cdot2048+16=4112$ samples are discarded, leaving 6128 post-transient samples per generated record.

\subsection{Nonlinear Sources and Added Noise}
The order-matched control and bounded surrogate are, respectively,
\begin{align}
 g_3(s)&=G s|s|^2,\label{eq:pure-cubic}\\
 g_{\tanh}(s)&=G\tanh(\gamma|s|)\frac{s}{|s|+\epsilon},\label{eq:tanh}
\end{align}
with $G=1.2+3.7j$, $\gamma=2$ and $\epsilon=10^{-12}$. Apart from the common phase rotation in $G$, the second mapping preserves input phase. For $|s|\gg\epsilon$, its local series is
\begin{equation}
 g_{\tanh}(s)=G\left(\gamma s-\frac{\gamma^3}{3}s|s|^2
 +\frac{2\gamma^5}{15}s|s|^4+\cdots\right).
 \label{eq:tanh-series}
\end{equation}
The leading small-signal response is linear, not cubic; receiver-band selection suppresses the linear component. Higher-order measurements and connector models motivate studying departures from a cubic law \cite{Christianson2008HigherOrder,Henrie2008Prediction}. However, this radial complex-baseband tanh mapping is a controlled surrogate, not a demonstrated exact baseband equivalent of a measured real-valued connector I--V law. Even the cubic source need not lie exactly in a dictionary with truncated memory $M<11$.

For noisy records, circular complex AWGN is added after the receiver filter: $y=p+v$, with $v=\sqrt{\sigma_v^2/2}(a_v+jb_v)$ and independent standard-normal components. Its variance is $\sigma_v^2=P_p10^{-L/10}$, where $P_p$ is clean PIM power over the 6128 post-transient samples and $L$ is the configured PIM-to-noise ratio. This simulation calibration does not enter model selection as a clean target. The historical condition is $L=35$ dB; the revised grid adds 20 and 50 dB and a clean condition.

\subsection{Cancellation Metric and Historical Split}
For an evaluation index set $\mathcal I$, cancellation depth is
\begin{equation}
 C_{\mathcal I}=10\log_{10}\frac{\operatorname{mean}_{n\in\mathcal I}|y[n]|^2}
 {\operatorname{mean}_{n\in\mathcal I}|y[n]-\widehat p[n]|^2+10^{-30}}.
 \label{eq:metric}
\end{equation}
Observed cancellation uses $y=p+v$; PIM-only cancellation replaces $y$ by $p$ and is a simulation diagnostic. Selection always uses observed validation cancellation. Neither $R^2$ nor an amplitude-spectrum ordinate substitutes for \eqref{eq:metric}.

The historical notebooks use TX seed 12345 and noise seed 54321. Their fitting and test intervals are $[4112,7176)$ and $[7176,10240)$, with 3064 samples each and no validation set. Several operating points were selected by maximizing the displayed test score. These values are exploratory and may be optimistically biased; relabeling that segment cannot repair the selection protocol.

\subsection{Independent Selection and Repetition Protocol}\label{sec:selection}
The revision generates three independent records for fitting, validation and test. To retain the original fitting count, each stage uses the first 3064 post-transient samples, $[4112,7176)$, of its own 10240-sample record; the remaining samples are not scored. Each record starts with its own filter transient. The counts are therefore $N_f=N_v=N_t=3064$, with no shared waveform history between stages. These are sample counts, not counts of statistically independent observations within a filtered record.

For repetition $r$ and stage $s=0,1,2$, TX and noise seeds are $b+1000r+10s$ and $b+1000r+10s+500000$. The fixed benchmark base is $b=30261001$; the infrastructure pilot uses a disjoint base $20261001$. All methods share the records in a repetition. Ten repetitions are specified for each source/dictionary and noise condition. Normalization and screening use fitting samples only. Parameters are chosen by the largest observed validation cancellation among converged finite solutions, with exact ties resolved by fewer active coefficients and then smaller $M$. Coefficients and selection metadata are frozen before generating the test record; no fitting-plus-validation retraining is performed. Nonconverged configurations and resource failures remain in the log.

The frozen common benchmark grid is $M\in\{2,4,6\}$. The common preflight estimate $16(8P^2+6N_fP)$ bytes is approximately 45.10 GiB for the combined $M=8$ dictionary, exceeding the 10 GiB workspace budget. Therefore $M=8$ is excluded for every method in the new comparison, including the cheaper cubic-only cases. Historical $M=8$ results are diagnostics and are not pooled with the revised benchmark.

\subsection{Ablations, Baselines and Numerical Settings}
The four ablations are full ridge-LS, conventional fixed-noise SBL without screening, screening followed by ridge-LS, and screening followed by SBL. Complex LASSO and complex orthogonal matching pursuit (OMP) provide additional sparse baselines. LASSO penalizes $\sum_i|w_i|$ in fitting-normalized coordinates, keeping the real and imaginary parts of one complex coefficient together; OMP also selects whole complex monomials. Both use the same candidate dictionary and receiver filter.

The ridge grid is $\rho=10^k$ for $k=-8,\ldots,-3$, with six values including the historical $10^{-4}$ reference. Other grids are $\beta\in\{10^2,10^4,10^6\}$ and $\eta\in\{0.01,0.65,0.8\}$. LASSO uses penalties $\{10^{-4},10^{-3},10^{-2},10^{-1}\}\lambda_{\max}$, where $\lambda_{\max}=\|\bar X_f^H\bar y_f\|_\infty$; OMP budgets are 5, 10, 20 and 40 terms. Each method uses the same validation observations, although grid cardinalities differ. Matched validation ablations compare screening and no screening at the same $M$ and regularization parameter; the four tuned pipelines are also compared on untouched test data.

For SBL, $\alpha_0=1$, $\alpha_{\max}=10^6$, denominator floor $10^{-12}$ and relative-mean tolerance $10^{-6}$ are fixed. The stopping ratio is $\|\mu_t-\mu_{t-1}\|_2/\max(\|\mu_{t-1}\|_2,10^{-12})$. The cap is 2000, versus 500 for the historical main sweep. Pruning resets the previous-mean comparison; at least one coefficient remains. LASSO uses complex coordinate descent, global relative KKT tolerance $10^{-5}$ and a 100000-sweep cap. The earlier infrastructure pilot used a different proximal solver and is excluded from the final comparison. Full trajectories are saved for selected SBL models; iteration counts, convergence flags and final stopping diagnostics are retained for every candidate.

The benchmark ran on an AMD Ryzen 5 2500U (four physical cores, 18.58 GiB RAM), Linux, Python 3.14.4, NumPy 2.5.3 and SciPy 1.18.1, using one BLAS thread per worker; build metadata are archived. Six one-hour timeouts were retried with a three-hour wall limit and unchanged scientific settings, producing all 720 results. Concurrent-search timings are retained for provenance but not used for method-speed comparisons.

Separate sequential profiles fix repetition 0, PNR 35 dB and all three source/dictionary cases before measurement. Each selected model is refitted once in a fresh process pinned to CPU 0, with one BLAS thread, and checked against its frozen coefficients. Dictionary construction is timed separately; fit time includes normalization, correlation computation, screening and solving. The implementation computes correlations even for unscreened branches. Inference includes active-monomial construction and receiver filtering over 3064 samples, with one warm-up and seven timed evaluations. Linux process high-water RSS includes imports, records, fitting and inference, but excludes the validation search. These are representative implementation measurements, not population estimates of runtime.

Means, standard deviations and paired confidence intervals use ten independent repetitions, not correlated samples within a record. Support stability uses monomial identifiers and pairwise Jaccard overlap within equal selected $M$. Pairs sharing a repetition are dependent, so only descriptive means and pair counts are reported; no confidence interval treats those pairs as independent. Inclusion frequencies are conditioned on selected memory. For a truncated tanh model this is stability, not recovery of a known exact finite support.

\section{Verified Diagnostics and Revision Results}\label{sec:results}
\subsection{What the Original Screen Retains}
Table~\ref{tab:screen-counts} reports the directly recomputed counts for all five historical scenarios. Screening is independent of $\beta$ and retains every column for every original configuration. Scores over all tested columns range from approximately 0.587 to 0.991. The previously quoted lower bound 0.61 omitted the extreme at combined $M=8$. These observations concern the supplied realizations; frequency-cell restriction alone does not prove that every possible target-correlation screen is ineffective.
\begin{table}[!t]
\caption{Historical candidate counts before/after screening at $\eta=0.01$. Counts apply to all original $\beta$ values and both clean/noisy versions where present.}\label{tab:screen-counts}
\centering\small
\begin{tabular}{rrr}\toprule
$M$ & Cubic: full/screened & Cubic--quintic: full/screened\\\midrule
2 & 6/6 & 40/40\\4 & 40/40 & 760/760\\
6 & 126/126 & 4788/4788\\8 & 288/288 & 18336/18336\\\bottomrule
\end{tabular}
\end{table}
Consequently, screen-plus-SBL and no-screen SBL have the same input dictionary at this threshold; the same holds for the two ridge branches. Final historical supports of seven cubic terms or 17 combined terms must be attributed to relevance pruning, not to a reduction by the first stage. Tuned higher thresholds require separate validation and may discard useful correlated columns.

\subsection{Iteration-Cap and Initialization Diagnostic}
For the original cubic source at $M=6$, $\beta=10^6$, raising the cap from 500 to 2000 changes cancellation from 51.264645 to 51.265265 dB. With $\alpha_0=1$, convergence occurs at iteration 1623; the last relative change in the posterior mean is $9.9881\times10^{-7}$, compared with $2.4041\times10^{-5}$ at iteration 500. Table~\ref{tab:convergence} reports initialization sensitivity. A cap of 5000 yields the same stopping points as 2000. All converged runs select the same seven monomials, and their cancellation differs by less than $4\times10^{-8}$ dB. Per-iteration residual power, negative log marginal likelihood and relative-mean changes are provided as machine-readable trajectories. This supports numerical stability on the historical record, not unbiased model selection or global optimality.
\begin{table}[!t]
\caption{Historical cubic $M=6$, $\beta=10^6$, cap 2000.}\label{tab:convergence}
\centering\small
\begin{tabular}{rrrr}\toprule
$\alpha_0$ & Iterations & Cancellation (dB) & Active\\\midrule
0.1 & 1603 & 51.265265 & 7\\
1 & 1623 & 51.265265 & 7\\
10 & 1612 & 51.265265 & 7\\\bottomrule
\end{tabular}
\end{table}

For $\alpha_0=1$, the normalized fitting residual power at iterations 1, 100, 500 and 1623 is respectively $0.586$, $2.114$, $2.141$ and $2.149$ in units of $10^{-9}$. The corresponding negative log marginal likelihood is $-42205.659$, $-42274.070$, $-42274.078331$ and $-42274.078448$ (additive Gaussian constant omitted). The residual need not decrease during pruning; the near-constant late objective and shrinking parameter change, rather than residual monotonicity, support the numerical stability statement.

\subsection{Why the Historical LS Margin Is Not a Final Comparison}
At the same historical cubic $M=6$, the saved ridge sensitivity study gives 50.682 dB for $\rho=10^{-4}$ but 57.891 dB for $\rho=10^{-8}$, against about 51.265 dB for SBL. The advertised margin over fixed ridge is therefore not evidence of superiority over a competitive tuned baseline. These alternative ridge values were also evaluated on the original test segment; they motivate validation tuning rather than supplying a replacement unbiased estimate. The historical tanh results, approximately 24 dB with a cubic dictionary and 30 dB with a combined dictionary, similarly remain exploratory controls.

\subsection{Independent Benchmark}
All 720 independent selection/test tasks completed, covering 15,840 candidate fits. Each cell in Table~\ref{tab:benchmark} summarizes ten independently generated fitting/validation/test triplets. The reported dispersion is the sample standard deviation across repetitions; means are computed in dB. The supplement also provides PIM-only scores, unadjusted 95\% Student-$t$ intervals and paired method differences.

\begin{table*}[!t]

\caption{Untouched-test observed cancellation (dB), mean $\pm$ SD over ten repetitions. All six methods use the same fitting/validation/test records and the frozen grids in Section~\ref{sec:selection}.}

\label{tab:benchmark}\centering\small\setlength{\tabcolsep}{3pt}

\begin{tabular}{llrrrr}\toprule

Source / dictionary & Method & Clean & PNR 20 dB & PNR 35 dB & PNR 50 dB\\\midrule

Cubic / cubic & Ridge-LS & $58.02\pm0.87$ & $20.12\pm0.54$ & $35.05\pm0.54$ & $49.41\pm0.54$\\

 & SBL & $51.96\pm0.65$ & $20.12\pm0.54$ & $35.00\pm0.53$ & $47.93\pm0.48$\\

 & Screen + ridge & $58.02\pm0.87$ & $20.12\pm0.54$ & $35.05\pm0.54$ & $49.42\pm0.54$\\

 & Screen + SBL & $51.96\pm0.65$ & $20.12\pm0.54$ & $35.00\pm0.53$ & $47.93\pm0.48$\\

 & Complex LASSO & $44.65\pm1.16$ & $20.12\pm0.54$ & $34.63\pm0.54$ & $43.53\pm0.95$\\

 & Complex OMP & $57.90\pm0.87$ & $20.12\pm0.54$ & $35.03\pm0.53$ & $49.40\pm0.54$\\

\midrule

Tanh / cubic & Ridge-LS & $24.75\pm1.22$ & $18.78\pm0.39$ & $24.34\pm1.07$ & $24.73\pm1.21$\\

 & SBL & $24.78\pm1.23$ & $18.79\pm0.41$ & $24.36\pm1.09$ & $24.76\pm1.22$\\

 & Screen + ridge & $24.76\pm1.22$ & $18.78\pm0.39$ & $24.35\pm1.08$ & $24.75\pm1.21$\\

 & Screen + SBL & $24.76\pm1.23$ & $18.80\pm0.40$ & $24.36\pm1.09$ & $24.74\pm1.22$\\

 & Complex LASSO & $24.81\pm1.19$ & $18.81\pm0.40$ & $24.39\pm1.07$ & $24.80\pm1.18$\\

 & Complex OMP & $24.68\pm1.19$ & $18.77\pm0.40$ & $24.29\pm1.10$ & $24.66\pm1.18$\\

\midrule

Tanh / cubic--quintic & Ridge-LS & $29.43\pm0.61$ & $19.53\pm0.42$ & $28.31\pm0.50$ & $29.38\pm0.60$\\

 & SBL & $29.32\pm0.69$ & $19.31\pm0.38$ & $28.29\pm0.47$ & $29.27\pm0.66$\\

 & Screen + ridge & $29.46\pm0.59$ & $19.58\pm0.46$ & $28.36\pm0.49$ & $29.42\pm0.58$\\

 & Screen + SBL & $29.42\pm0.61$ & $19.47\pm0.45$ & $28.33\pm0.46$ & $29.38\pm0.62$\\

 & Complex LASSO & $29.40\pm0.66$ & $19.57\pm0.45$ & $28.32\pm0.56$ & $29.35\pm0.65$\\

 & Complex OMP & $28.95\pm0.91$ & $19.51\pm0.43$ & $27.99\pm0.59$ & $28.98\pm0.77$\\

\bottomrule\end{tabular}\end{table*}

\begin{table}[!t]

\caption{Mean dictionary sizes at PNR 35 dB over ten independently selected models. Fractional counts reflect different selections across repetitions. Per-configuration counts for all 15,840 fits are supplied in \texttt{candidate\_results.csv}.}

\label{tab:stage-counts}\centering\scriptsize

\begin{tabular}{lrrr}\toprule

Method & $N_{\rm full}$ & $N_{\rm screen}$ & $N_{\rm active}$\\\midrule

\multicolumn{4}{l}{\emph{Cubic / cubic}}\\

Ridge-LS & 126.0 & 126.0 & 126.0\\

SBL & 126.0 & 126.0 & 7.8\\

Screen + ridge & 126.0 & 113.2 & 113.2\\

Screen + SBL & 126.0 & 119.3 & 7.9\\

Complex LASSO & 126.0 & 126.0 & 12.2\\

Complex OMP & 117.4 & 117.4 & 18.0\\

\midrule

\multicolumn{4}{l}{\emph{Tanh / cubic}}\\

Ridge-LS & 126.0 & 126.0 & 126.0\\

SBL & 100.2 & 100.2 & 24.3\\

Screen + ridge & 117.4 & 108.3 & 108.3\\

Screen + SBL & 100.2 & 94.8 & 22.4\\

Complex LASSO & 126.0 & 126.0 & 12.5\\

Complex OMP & 117.4 & 117.4 & 11.5\\

\midrule

\multicolumn{4}{l}{\emph{Tanh / cubic--quintic}}\\

Ridge-LS & 4788.0 & 4788.0 & 4788.0\\

SBL & 3176.8 & 3176.8 & 30.0\\

Screen + ridge & 4788.0 & 2892.0 & 2892.0\\

Screen + SBL & 4385.2 & 2396.9 & 24.0\\

Complex LASSO & 4788.0 & 4788.0 & 33.1\\

Complex OMP & 4788.0 & 4788.0 & 20.0\\

\bottomrule\end{tabular}\end{table}

For the clean cubic source, tuned ridge reaches 58.02 dB, compared with 51.96 dB for SBL and 57.90 dB for OMP. Thus the revised comparison does not support an unconditional cancellation advantage for SBL. At PNR 35 dB, ridge and SBL give 35.05 and 35.00 dB, while SBL uses 7.8 active coefficients on average instead of 126.

For tanh at PNR 35 dB, increasing the dictionary from cubic to cubic--quintic raises ridge cancellation from 24.34 to 28.31 dB. In the combined dictionary, screen-plus-SBL gives 28.33 dB with 24.0 active coefficients on average, compared with 28.31 dB and 4788 coefficients for ridge. The selected memories and thresholds can differ between methods; these are comparisons of validation-tuned pipelines.


\subsection{Screening Ablation and Numerical Stability}
For every configuration, the supplement records $N_{\rm full}$, $N_{\rm screen}$, $N_{\rm active}$ and convergence. At $\eta=0.01$, all 3240 matched ridge/SBL candidate pairs have identical validation scores, counts and convergence status with and without screening. Thus this original threshold contributes no reduction. Across the three thresholds there are 9720 matched validation controls; these isolate screening at fixed memory and regularization without retuning on test data. The four validation-tuned test branches in Table~\ref{tab:benchmark} address the end-to-end tradeoff and should not be interpreted as fixed-parameter causal effects.

For tanh/cubic--quintic at PNR 35 dB, the paired test difference of screen-plus-SBL minus SBL is 0.039 dB, with an unadjusted 95\% interval of $[-0.027,0.106]$ dB over ten repetitions. This does not establish a cancellation gain from the tuned screen. These intervals are descriptive and are not corrected for the multiple method/condition comparisons.

Of 15,840 candidates, 536 did not satisfy the stopping criterion: 238/1440 LASSO fits, 86/1080 SBL fits and 212/3240 screen-plus-SBL fits. All were excluded from selection, and all 720 selected models converged; no repetition was dropped. The 240 selected SBL trajectories and historical cap/initialization controls are archived. This avoids using a capped solution as the main comparison, while making the finite optimization budget explicit.

At PNR 35 dB the cubic SBL support has mean Jaccard overlap 0.364 across 45 pairs of repetitions. For tanh/cubic--quintic, the corresponding values are 0.147 over 21 equal-memory pairs for SBL and 0.095 over 36 pairs for screen-plus-SBL. These overlaps are descriptive and conditioned on selected memory; dense ridge's overlap of one simply reflects retention of the whole dictionary. The much lower sparse-support overlaps, despite similar cancellation, preclude interpreting the selected monomials as uniquely recovered physical components.

\subsection{Resource Measurements}
\begin{table*}[!t]
\caption{Sequential resource measurements: one refit of the frozen selected model per method/case, repetition 0 at PNR 35 dB. Times exclude hyperparameter search. Inference is the median of seven warmed block evaluations, including active-feature construction and receiver filtering. Peak is the process high-water RSS.}
\label{tab:resources}\centering\scriptsize\setlength{\tabcolsep}{3pt}
\begin{tabular}{llrrrrrr}\toprule
Source / dictionary & Method & $M$ & Active & Dictionary (s) & Fit (s) & Peak (MiB) & Inference ($\mu$s/sample)\\\midrule
Cubic / cubic & Ridge-LS & 6 & 126 & 0.070 & 0.052 & 151.8 & 21.27\\
 & SBL & 6 & 8 & 0.068 & 0.082 & 138.1 & 0.80\\
 & Screen + ridge & 6 & 114 & 0.072 & 0.049 & 146.1 & 20.54\\
 & Screen + SBL & 6 & 8 & 0.070 & 0.081 & 138.0 & 0.80\\
 & Complex LASSO & 6 & 15 & 0.069 & 2.505 & 138.3 & 1.67\\
 & Complex OMP & 6 & 20 & 0.071 & 0.080 & 138.0 & 2.17\\
\midrule
Tanh / cubic & Ridge-LS & 6 & 126 & 0.070 & 0.053 & 151.5 & 21.44\\
 & SBL & 4 & 39 & 0.022 & 0.063 & 128.9 & 6.07\\
 & Screen + ridge & 6 & 114 & 0.070 & 0.046 & 146.2 & 19.58\\
 & Screen + SBL & 4 & 39 & 0.022 & 0.062 & 128.5 & 5.92\\
 & Complex LASSO & 6 & 14 & 0.068 & 3.011 & 137.9 & 1.50\\
 & Complex OMP & 6 & 10 & 0.070 & 0.043 & 138.1 & 1.02\\
\midrule
Tanh / cubic--quintic & Ridge-LS & 6 & 4788 & 2.642 & 15.646 & 1462.1 & 891.53\\
 & SBL & 6 & 19 & 2.675 & 243.413 & 3439.4 & 2.11\\
 & Screen + ridge & 6 & 2011 & 2.758 & 5.129 & 889.1 & 361.77\\
 & Screen + SBL & 6 & 18 & 2.805 & 23.471 & 1124.7 & 1.98\\
 & Complex LASSO & 6 & 33 & 2.940 & 109.578 & 1377.4 & 7.21\\
 & Complex OMP & 6 & 20 & 2.750 & 1.882 & 1014.1 & 2.24\\
\bottomrule\end{tabular}\end{table*}

Table~\ref{tab:resources} separates fitting cost from deployment cost. In the representative combined case, SBL and screen-plus-SBL take 243.41 and 23.47 s to fit, with process peaks of 3439 and 1125 MiB; ridge takes 15.65 s. Screen-plus-SBL therefore remains slower to fit than ridge in this slice, while its sparse inference takes 1.98 versus 891.53 $\mu$s per scored sample. All 18 profiled models reproduce their frozen coefficients exactly. Different selected memories and screened dictionaries limit these comparisons to the declared slice. Building the full dictionary remains necessary before screening; a small final support does not imply an inexpensive fit. The seven inference repetitions describe execution variability for that block, not variability across independently fitted models. No speed advantage is inferred from concurrent full-search timings.

\section{Discussion and Limitations}\label{sec:discussion}
The data distinguish order matching from order mismatch: a cubic source has a matching interaction family but still suffers finite-memory truncation, whereas tanh introduces higher odd orders. Sparse supports can approximate a dense fit without identifying unique physical coefficients because narrowband Volterra columns are strongly correlated. Stability of cancellation and stability of selected monomials must therefore be evaluated separately.

The frozen grids are a material limitation. LASSO chooses its smallest penalty in 118/120 tasks, unscreened SBL chooses the largest $\beta$ in 102/120, and ridge reaches its lower or upper regularization boundary in 23/120 and 46/120 tasks. The common maximum $M=6$ is also often selected. These validation-selection diagnostics show that the grids do not establish globally competitive optima for every estimator. We retain the predeclared comparison and report boundary frequencies instead of enlarging ranges after inspecting test results; wider-grid ranking claims require a new evaluation. The conclusions concern compactness and screening behavior under this specified budget, not an intrinsic superiority of SBL over LASSO or ridge.

Compactness alone does not establish computational efficiency. Even a final seven-term model can require building a full dictionary and repeatedly solving a large system. At combined $M=8$, one complex128 $18336\times18336$ Gram matrix occupies about 5.01 GiB, before other arrays. The present implementation and resource-limited benchmark must not inherit speedup claims from an untested incremental variant. No residual-ranked incremental method is evaluated in this revision.

The posterior equations describe regularized coefficient inference. Since the likelihood precision is externally tuned, the tanh source is misspecified by a finite-order dictionary and no predictive-coverage analysis is reported, calibrated predictive uncertainty is not claimed. Real measured PIM, changing paths, multiple nonlinear sources and MIMO remain outside this synthetic study \cite{Lampu2023Algorithms,Lampu2023Cascade}. Testing measured data is future work.

\section{Conclusion}\label{sec:concl}
The repeated validation-selected benchmark demonstrates a compactness--accuracy tradeoff for fixed-noise SBL in synthetic PIM cancellation. Tuned ridge outperforms SBL on the clean cubic source, whereas at 35 dB PIM-to-noise ratio both approach the noise floor with markedly different coefficient counts. For the tanh source, fifth-order terms improve cancellation, and screen-plus-SBL retains a small fraction of the combined dictionary. The original $\eta=0.01$ screen is inactive; higher-threshold screening and relevance pruning have distinct effects. Sequential resource measurements and low sparse-support overlaps limit claims of computational efficiency and physical identifiability. Wider tuning ranges, measured PIM and changing channels remain necessary before practical superiority can be established.

\bibliographystyle{IEEEtran}
\IEEEtriggeratref{10}
\bibliography{references}
\end{document}
